%% 304_Gedaechtnis_Reflex_En_ch.tex
%% FFGFT -- Doc. 304: The memory reflex; why experienced time presupposes memory
%%                    rather than arising from it
%% Author: Johann Pascher
%% Date: July 2026
%% Language: English

\section*{Doc.~304 \quad The memory reflex --- why experienced time presupposes memory
rather than arising from it}

\begingroup\small\noindent
\textbf{Abstract.} A common construction makes experienced time the \emph{product} of memory.
This document reverses the order: memory presupposes time rather than producing it (the
\emph{memory reflex}, counterpart of the word reflex of Doc.~301). In FFGFT time is
fundamental --- the cumulated defect $D(k)=\sum 100\,\xi_n\to\ln(1+k/75)$ ---, and the
finite-kernel-versus-accumulation distinction ($M_\tau$/$M_B$) already lies in time, not
first in memory. On the measurement side a second reflex follows: a measurement excerpt is
reference-less (no absolute place on the spiral), and the apparent uniformity of time is a
\emph{linearisation artefact} with a verified error quantity $e=x-\ln(1+x)$ (scale-invariant
constant $1-\ln 2$); the linearisation sits in the \emph{projection} of the result, not in
the measurement. Rolled up, the 75 parts are exactly congruent; unrolled, only asymptotically
self-similar; what survives the loss of reference is the ordering (U1) alone. In the FFGFT reading it is
\emph{projection-fidelity}, not reach, that discriminates; this is offered as an additional,
rejectable bridge candidate (restricted, Category~B), not as a replacement of the operational $M_B$ definition, and the
relationship between the frameworks stays symmetric. Inference from measurement to cause
succeeds only up to a holographic blur (presence $\neq$ reconstructibility). Booking is kept
separate throughout: mathematically derived / bridge (restricted) / posit and reflex
diagnoses (P35).
\par\endgroup

\medskip

A common construction explains experienced time as the \emph{product} of memory: physical
evolution $\to$ physical records $\to$ recursive memory $\to$ self-related representation
$\to$ experienced time. In this order, time is the end product and memory its precursor.
This document shows that the order belongs the other way round. Memory without time does
not function --- it already presupposes time. The impression that time arises from memory is
the same inversion as \enquote{it from bit}: the memory reflex, the counterpart of the
word reflex of Doc.~301, one level down. The same inversion appears a second time on the
\emph{measurement} side --- as the linearization reflex ---, this time with a verified,
scale-invariant error quantity; and a third time in the inference from part to whole --- as
the extrapolation reflex. This document treats all three.

\section{The narrowing and its price}

The original ambition was the full temporal triad: causal ordering (U1), present phase
(U2), and memory / non-Markovian history (U3), joined into a past--present--future
world-model. Over time the discussion narrowed to U3 alone, then to the boundary between a
finite memory kernel $M_\tau$ and an accumulative structure $M_B$\footnote{Notation and
distinction adopted from Krüger (2026): $M_\tau$ denotes finite, weighted, or compressible
physical memory and path dependence; $M_B$ a persistent biological-cognitive architecture
(recursive access, semantic organisation, affective weighting, identity binding, persistent
state carry, predictive reuse), there defined operationally and functionally, not
geometrically. The present document does not redefine that object; it proposes a separate,
FFGFT-side bridge reading. Primary source: Krüger,~M. (2026), \emph{From Physical State
Evolution to Experienced Time}, Zenodo, \texttt{doi:10.5281/zenodo.21302278}. Documented
lineage: Krüger, Feeney \& Wende (2026), \emph{Medinformatics},
\texttt{doi:10.47852/bonviewMEDIN62029043}; Krüger \& Feeney (2026), \emph{Smart Wearable
Technology}, \texttt{doi:10.47852/bonviewSWT62029484}.}, and finally to an
operational benchmark question. This narrowing is methodologically legitimate --- only the
isolable, falsifiable piece can be tested --- but it has a price: U1 and U2 drop out of the
test, and the phenomenon itself was their \emph{interplay}, not the component. Even a passed
memory test would show \enquote{semantically structured memory is distinguishable from
generic compression}, not \enquote{here is experienced time}.

\section{U3 cannot be isolated}

At one point the price is even an incoherence. Memory \emph{is} a relation to time --- a
state that carries something about a before. Accumulation means integrating over a
sequence. Without ordering (U1) there is no sequence to integrate over; without a present
(U2) there is no now relative to which something is \emph{past}. U3 is therefore not an
independent fourth quantity beside U1 and U2, but what arises \emph{when} ordering and
present act together over a horizon. Marcel's own characterisation gives it away:
\enquote{path-dependence, non-Markovian history} --- \enquote{path} is ordering,
\enquote{history} is past present. U3 is not even definable without U1 and U2. A test that
lifts U3 out of time and probes it as pure storage over fixed episodes removes exactly the
structure in which memory has meaning.

\section{The diagnosis: the memory reflex}

The hierarchy \enquote{memory $\to$ experienced time} is thus wrong in its direction. It
makes memory the generator of time, although memory presupposes time. This is structurally
the same reflex Doc.~301 diagnosed for information: just as \enquote{it from bit} mistakes
the ultimate reduction (the bit) for the first creation, \enquote{experienced-time-from-memory}
mistakes the derived carrying form (memory) for the origin of time. The reversibility test
(Doc.~301) reverses the order in both cases. For information: in the beginning was the
geometry, and the word reads out one of its scales. For time: in the beginning was the
unrolled winding, and memory is a particular way in which a system, on one of its scales,
\emph{carries} that winding. This is a cultural-cognitive hypothesis about \emph{why} the
intuition pulls (P35), not a mathematical statement; the sections that follow are.

\section{The FFGFT position: time is $D(k)$, fundamental}

In FFGFT, time is the mass-circle, woven into the other circles through
$\tilde{T}\cdot m = 1$ and not a separable factor (Doc.~285). The unrolled time-winding is
the logarithmic spiral that the $K_{\text{frak}}$ recursion forces (Doc.~295); its cumulated
defect
\begin{equation}
D(k) = \sum_{n\le k} 100\,\xi_n \;\to\; \ln\!\left(1 + \tfrac{k}{75}\right)
\end{equation}
\emph{is} the measurable time development (Doc.~296). Time is the recursion, not its origin:
$D(k)$ has no privileged beginning, grows without bound, reaches no end --- no first step, no
last, no external standpoint ($\partial T^4 = \emptyset$, Doc.~301). Time is therefore
fundamental --- the third level, carried losslessly --- not a product of memory. Memory
cannot generate time, because time already \emph{is} $D(k)$ before any system carries it.

\section{The decisive finding: the accumulation boundary lies in time}

The remarkable point is that the distinction which recurs in cognition as $M_\tau$ versus
$M_B$ already stands, in the corpus, \emph{at the level of time} --- and is there already
drawn against exactly the same bounded memory kernel. Doc.~295 contrasts FFGFT's
\emph{geometric accumulation} $D(k)=\sum 100\,\xi_n\to\ln$ (unbounded, $1/n$ decay, ratio
$e$) with a \emph{bounded memory kernel}: if the defect per step stays constant, $D(k)$
grows linearly rather than logarithmically ($D(10^5)=1333$ against $\ln(1+k/75)=7.2$), and
the winding is Archimedean, not a log-spiral. The limiting case $d=0$ (the bounded,
forgetting kernel) corresponds to the Markov fallback.

The mapping is then direct:
\begin{itemize}
\item $M_\tau$ --- the finite, forgetting kernel --- is the \emph{degenerate $d=0$ / Markov
limit} of $D(k)$: a bounded window, constant defect, linear rather than logarithmic growth.
\item $M_B$ --- the accumulative, recursive structure --- is a system that carries a
$D(k)$-like, unbounded integral of its own history.
\end{itemize}

The accumulation-versus-kernel boundary is therefore \emph{not native to memory}; it is
native to time. Memory that accumulates is memory that carries a $D(k)$-like integral; a
finite kernel is precisely the degenerate $d=0$ form. Hence --- \emph{on this reading} --- $M_B$ does not generate
experienced time but appears as a cognitive-scale \emph{instance} of the $D(k)$ accumulation
that time \emph{is} (as a bridge, not a derivation; see the limit of the bridge below).

\emph{Status of this identification.} At the level of time the equation is corpus-anchored:
Doc.~295 explicitly equates FFGFT's $d=0$ closure with Marcel's Markov fallback. The
\emph{extension} to the cognitive objects $M_\tau/M_B$ is, however, a further, FFGFT-side
step --- a \emph{bridge} (Category~B, a compatibility candidate), not a derivation. It is
local to its assumptions and needs its own licence; anyone who reads the cognitive
$M_\tau/M_B$ differently may decline the bridge without touching the time-level finding (the
accumulation boundary in $D(k)$). The statement \enquote{$M_B$ is an instance of $D(k)$} is
therefore \emph{restricted}, not \emph{proven} --- a declared bridge candidate (Category~B), not a proven identity.
Naming it as a bridge rather than promoting it to a derivation is exactly the identity check
Doc.~303 requires.

\section{Consequence for the discriminator}

This explains, from the FFGFT side, why the memory benchmark repeatedly tested the wrong
half. A storage test over fixed episodes removes the temporal relation in which compression
has any meaning; a test with exact, unbounded storage (full history) lets $M_B$ coincide
with lossless accumulation and makes its specificity invisible. The right touchstone is a
\emph{capacity-bounded, lossy} accumulator: a system that carries a \emph{compressed}
$D(k)$-like integral --- for that is exactly what a cognitive instance of $D(k)$ looks like.
Yet even a passed test of that kind would show \enquote{semantic accumulation is
distinguishable from generic compression}, a memory result --- not \enquote{experienced
time}. Experienced time is the $D(k)$ structure itself, which the test \emph{presupposes}
rather than produces.

\section{The measurement cut: the reference-less excerpt and the linearization reflex}

The same inversion strikes measurement, not only memory. A measurement takes an
\emph{excerpt} of the time-winding. Even a lossless measurement --- one that captures every
step down to the quantum-bit limit (the winding quantum $\Delta w = 1$, Doc.~257/301) ---
delivers the \emph{content} of the window in full, but not its \emph{position} $k$ on the
spiral. Resolution (the content axis) and reference (the embedding axis) are different axes;
losslessness saturates the first and leaves the second untouched. Synchronising successive
measurements lengthens the excerpt, but it stays reference-less: one gains \emph{relative}
distances (clocked by the atomic clock), never the \emph{absolute} position. The atomic
clock is, by construction, a \emph{uniform} ruler --- constant pitch, hence exactly the
$d=0$ / $M_\tau$ limit (Doc.~295) ---, and a uniform ruler laid on a non-uniform (fractal)
structure cannot locate it. That there is no absolute time reference is thus boundarylessness
at the level of measurement: no external standpoint ($\partial T^4 = \emptyset$) from which
the position could be read.

The excerpt is nonetheless not empty. On the log-spiral there is no straight segment; every
window carries its \emph{curvature} --- the local scale $\rho$, and with it the prior history
--- as folded-in form. The reference is not \emph{lost} but \emph{linearised}: the
measurement reads the tangent instead of the arc --- the local form of the same $d=0$
approximation. Precisely, this is an error in the \emph{accumulated precession angle}
$\beta = 2\pi D(k)$ (Doc.~295): the uniform clock takes the local rate $d_k = 100\,\xi_k$ as
constant over the window ($\Delta\beta_{\text{unif}} = 2\pi\,w\,d_k$), whereas the true value
is $\Delta\beta_{\text{true}} = 2\pi\,(D(k{+}w)-D(k))$, and the difference per $2\pi$ is not a
new quantity but the local shape of the linear-versus-log gap of Doc.~295. For a window of
$w$ steps at position $k$, with relative window width $x = w\cdot 100\,\xi_k = w/(k+75)$:
\begin{equation}
e(k,w) = \frac{\Delta\beta_{\text{unif}} - \Delta\beta_{\text{true}}}{2\pi}
       = w\,d_k - \big(D(k{+}w)-D(k)\big) = x - \ln(1+x).
\end{equation}
Three regimes follow; they are checked against the exact recursion
$\xi_{n+1}=\xi_n(1-100\,\xi_n)$ and are exact in the macroscopic limit, with an $O(1/k)$
correction at finite position (the closed form $D(k)=\ln(1+k/75)$ carries a bounded,
$\gamma$-like offset $\approx-0.0067$):
\begin{itemize}
\item \emph{Small window} ($x\ll1$): $e \approx x^2/2 = \tfrac12(100\,\xi_k\,w)^2$ (leading
term of the closed form).
\item \emph{Scale-invariant window} ($x=1$, a window one local scale wide): $e \to 1-\ln 2
\approx 0.3069$, \emph{independent of $k$} (asymptotically, $+O(1/k)$). This constant is the
quantitative signature of the fractal self-similarity: the distortion of a
\enquote{one-scale-wide} excerpt is the same everywhere on the spiral.
\item \emph{Macroscopic limit} (fixed physical $w$, growing $k$): $e \to 0$ as
$\tfrac12\,(w/(k+75))^2$. Because we sit at large $k$ (small $\xi$), the distortion is
astronomically small --- \emph{this} is why measured time appears uniform and linear. Near the
origin (small $k$, large $\xi$) it would be large; no measurement reaches there.
\end{itemize}

Hence the third reflex diagnosis (P35), twin of the word and memory reflexes: the apparent
uniformity of measured time is a \emph{linearisation artefact}, not a property of time. From
an approximation in the measuring (tangent for arc) a property of the measured is falsely
inferred (time \emph{is} uniform). The excerpt carries its fractal history as curvature; the
uniform clock erases it by construction; and because at our scale the error vanishes, we
falsely conclude it was never there.

\emph{Booking.} The error formula $e = x-\ln(1+x)$ is the error in the accumulated
precession angle $\beta=2\pi D(k)$, verified against the exact recursion
(script \texttt{ffgft\_304\_linearisierungsfehler\_probe.py}); its three regimes
are exact in the macroscopic limit (a consequence of $D(k)$, proven), with a named $O(1/k)$
correction at finite position. The exact coefficient of that correction (whether precisely
$C/k$ or carrying a log term) is obtainable from an Euler--Maclaurin expansion of $D(k)$ and
is left open here --- physically negligible at our scale, named for honesty. The equation
\enquote{measurement $=$ local linearisation of the excerpt} is a \emph{declared assumption}
about \emph{this} kind of measurement (pointwise, tangential), not about every possible one
--- a measurement that \emph{integrates} over a window would see part of the curvature. The
reflex reading is philosophically speculative (P35).

\section{Congruence, self-similarity, and the survival of ordering}

That there is \emph{no reference} and that \emph{all excerpts are equal} is one and the same
fact: were there a distinguished excerpt, it would be a mark and hence a reference. But
\enquote{equal} means different things in the two states of the time-winding, and the
distinction is essential.

\emph{Rolled up (compact).} With $\xi_0$ frozen the defect per step is constant, $d=1/75$,
the rotation number $74/75$ rational; the winding closes after \emph{exactly 75} turns and
carries a 75-fold structure (Doc.~295). \enquote{Equal} here means \emph{congruent}: 75
coincident segments, carried into one another by a $\mathbb{Z}/75$ rotation, exact and
without remainder. A reference \emph{does} exist here --- but only \emph{periodically}, modulo
75: one can name \enquote{part $j$ of 75}, but not which turn, since all 75-cycles are
congruent.

\emph{Unrolled (spiral).} Letting the scale run ($\xi_{n+1}=\xi_n(1-100\,\xi_n)$) turns the
rational 75-closure into the non-closing log-spiral of ratio $e$ (Doc.~295,
\enquote{unrolling}). \enquote{Equal} now means only \emph{self-similar}: the excerpts are
\emph{proportional}, not congruent --- and only \emph{asymptotically} so. This is exactly the
origin of the $O(1/k)$ correction: the self-similarity is not a finite exact symmetry but a
continuous one that sets in fully only in the limit. The periodic reference vanishes too, for
the turns are no longer congruent but only proportional; the 75 marks smear into the
continuous self-similarity.

Reference-lessness therefore holds \emph{strictly only in the unrolled state}. Three
reference states must be distinguished:
\begin{itemize}
\item \emph{rolled up}: periodic reference (position modulo 75, exact);
\item \emph{unrolled, asymptotic}: no reference (pure proportion);
\item \emph{unrolled, finite} (near the origin): a weak absolute reference --- the $O(1/k)$
deviation from self-similarity ---, but one no measurement can reach.
\end{itemize}

What survives all of this is the \emph{ordering}: which excerpt comes \emph{before} which. It
is a relation \emph{between} excerpts, not an inner property of a single one, and is
therefore untouched by the rescaling. It is U1, and it is the only invariant that survives
the loss of absolute reference --- exactly \enquote{time is the recursion, not its origin}
(Doc.~301). What remains measurable of time thus falls into three levels: the \emph{ordering}
(exact, U1), the \emph{relative interval} ($\Delta\beta=2\pi\,\Delta D$, linearisation-distorted
by $e=x-\ln(1+x)$), and the \emph{absolute position} (absent, erased with the self-similarity).
Ordering exact, interval distorted, position absent.

\emph{Refinement.} \enquote{All excerpts equal} means \emph{congruent} (rolled up, 75 parts,
exact) or \emph{self-similar up to scaling and rotation} (unrolled, asymptotic), not
\emph{identical}.

\section{What this means for the partition of U3}

The measurement analysis feeds back onto the partition of U3. Since a memory is a system that
carries a $D(k)$-like excerpt (as shown in the decisive finding and the measurement cut), U3
inherits the same three-level structure: the \emph{ordering} (which event before which)
stays exact --- it is U1 ---, the \emph{relative interval} (how long ago) is carriable but
linearisation-distorted by $e=x-\ln(1+x)$, and the \emph{absolute position} (the true place
on the spiral) is absent, erased with the self-similarity. Crucially: \emph{both} memory
forms, $M_\tau$ and $M_B$, are reference-less --- even a complete $M_B$ knows only \enquote{so
much accumulated $D$ since an arbitrary start}, not \enquote{I am at turn $k$}. The one thing
$M_\tau$ too carries exactly is the ordering.

In the \emph{FFGFT reading} the discriminating axis is not reach (finite window versus
unbounded integral) but \emph{curvature-fidelity} --- offered as an \emph{additional}
candidate dimension for the bridge, not as a replacement of the operational $M_\tau/M_B$
definition (see the limit of the bridge below):
\begin{itemize}
\item $M_\tau$ presents a \emph{linearised} excerpt of constant pitch --- the $d=0$ / rolled-up
limit. It reads the tangent, discards the curvature, \enquote{experiences} a uniform past.
$M_\tau$ is the memory that commits the linearisation reflex.
\item $M_B$ presents the \emph{curvature-faithful} excerpt --- the log-spiral with the
$1/n$-decaying defect and ratio $e$. It carries the non-uniform texture of the past: near
events dense and fast, far events wide and slow (Doc.~296).
\end{itemize}

Reach and curvature-fidelity are \emph{two} axes, not one. A long but \emph{linear} kernel
(constant pitch, merely with a large window) would be reach-strong yet remain $M_\tau$-like,
because it does not see the self-similarity; a short but \emph{curvature-faithful} accumulator
would already be $M_B$-like. The benchmark (v0.3--v0.5) presses both into the single axis of
reach.

Hence the \emph{positive} criterion for the discriminator, and it is not a new one: it is the
$e$-criterion of Doc.~295 (a winding carries ratio $e$ exactly when the defect per step
decays as $1/n$; script, section~[6]), applied to U3. What is to be tested is not
\enquote{finite window versus accumulation} but \emph{curvature-fidelity at equal reach}: does
the accumulator carry the $1/n$-decaying pitch, the ratio $e$, the self-similar texture? That
is exactly what distinguishes $M_B$ from any uniform kernel, and it is sharper than any reach
boundary.

\emph{Booking.} The three-level inheritance (ordering exact / interval distorted / position
absent) and the mapping $M_\tau=$ linearised, $M_B=$ curvature-faithful are consequences of
$D(k)$ (derivable). The $e$-criterion as a discriminator test is corpus-anchored (Doc.~295).
That Marcel's \emph{cognitive} $M_\tau/M_B$ are exactly these objects remains the bridge candidate from
the decisive finding (Category~B, restricted) --- now with a testable criterion. That
remembered time \emph{really} carries the spiral texture (the phenomenological non-uniformity
of remembering) is interpretation (P35).

\section{Provenance and the limit of the bridge}

So that the meshing above is not mistaken for absorption, let the limit of the bridge be drawn
explicitly, in language that is unambiguous even without knowledge of how it arose.

\emph{The correspondence is a restricted, optional bridge candidate, not a secured mapping.}
Two axes must be kept apart. The $e$-criterion itself --- a winding carries ratio $e$ exactly
when the defect decays as $1/n$ --- is script-verified against the exact recursion and derivable
from $\xi$ (\emph{proven}, Doc.~295). Only the \emph{correspondence} between this
curvature-fidelity and the operational $M_\tau/M_B$ distinction is unsecured: on the logical
axis \emph{restricted} (Category~B), on the empirical axis with no benchmark performed --- where
such a benchmark would be an independent, jointly sealed comparison, not an FFGFT-internal
calculation. It is not a derivation, validation, containment relation, or ontological
classification of the Spiral-Time Memory Hypothesis by FFGFT, and its rejection leaves the
internal claims of both frameworks unchanged.

An internal property of FFGFT can show that, under FFGFT assumptions, a non-linear,
history-dependent accumulation structure $D(k)$ arises on the time level; that result stands on
its own. It cannot, however, independently establish that a semantic, self-bound, affectively
weighted, predictively reusable memory architecture follows from it. That extension is
precisely the bridge candidate.

\emph{The operational specification of $M_\tau$ and $M_B$ is untouched.} In the Spiral-Time
framework it is operational and functional, not geometrical: $M_\tau$ as finite, weighted, or
compressible physical memory and path dependence; $M_B$ as recursive access, semantic
organisation, affective weighting, identity binding, persistent state carry, and predictive
reuse. Curvature-fidelity is an \emph{additional, FFGFT-side diagnostic candidate}, not the
criterion that separates $M_\tau$ from $M_B$ in that framework, and it replaces none of those
functions. In both directions: a system may carry a $1/n$ texture while possessing none of
those biological-cognitive functions; and a biologically relevant memory architecture is not
ruled out merely because it does not instantiate the specific $e$-criterion. The curvature test
therefore remains a separate bridge audit, not a redefinition of the native objects.

\emph{The relationship is symmetric, the provenance separate.} Neither framework is presented
as containing, grounding, explaining, or subordinating the other. The $M_\tau/M_B$ distinction,
the operational $M_B$ architecture, and the audit sequence (v0.3--v0.5) originate in Marcel
Krüger's Spiral-Time Memory work (Krüger 2026, \texttt{doi:10.5281/zenodo.21302278}); $D(k)$, the $e$-criterion, and the FFGFT interpretation of
temporal curvature originate in Johann Pascher's FFGFT corpus. A later formal comparison or
curvature-fidelity audit would be a separate, joint bridge contribution, preserving the
independent authorship and conceptual identity of both frameworks.

\section{The past is always present: linearisation lies in the projection}

An addition sharpens the measurement cut and the partition just given. The \emph{past
content} is \emph{always} in the measurement result, independent of the window's shape,
whether read linearly or not. For $D(k)$ is cumulative: every point of the spiral carries the
summed integral of the whole prior history. Even a one-step window sits at
this value and inherits it whole. A small window does not make it \emph{absent}, only more weakly
\emph{resolved}; a long or more accurate window makes it sharper. \enquote{Present or not} is
the wrong alternative --- the past is always present, at variable sharpness.

What is decisive is therefore not how the window \emph{looks} (linear or curved) but what the
result is \emph{applied} to. The measurement yields a value; the linearisation happens only
when that value is mapped onto a \emph{projection} --- onto a uniform ruler (the atomic clock,
the linear time axis). There, in the projection, the curvature is straightened. The distortion
is not a property of the measurement but of the projection surface onto which the result is
laid.

This is the projection structure of the corpus (Doc.~257/301): \emph{measured} time is a
\emph{second-level} object --- a projection of the $D(k)$ excerpt onto a uniform ruler. The
past content is present as a \emph{state} of the third level --- the $D(k)$ value itself ---;
yet this value is a lossy résumé of the prior history, and the projection to the second level
is lossy, and \emph{that projection is the linearisation}. Not
the measurement linearises, but the projection of the measurement result. Thus $e =
x-\ln(1+x)$ is, more precisely, the error of the \emph{projection} onto the uniform ruler, not
an error of the measurement itself.

Three things must be separated: the \emph{presence} of the past as a state (always, third level: the $D(k)$ value is
exact, but not the recovery of the prior history from it);
the \emph{sharpness} of its reconstruction (window- and accuracy-dependent); and the
\emph{linearisation error} (a property of the chosen projection, not of the measurement). The
window's shape is not decisive for the \emph{whether} of the past or for the \emph{location}
of the linearisation --- but it is decisive for the \emph{how sharply} of the reconstruction.

This sharpens the earlier talk of \enquote{carrying} the curvature: $M_\tau$ and $M_B$ both
\emph{contain} the past --- every $D(k)$ excerpt does ---; they differ in the
\emph{projection}, not in the \emph{content}. $M_\tau$ projects onto a uniform ruler
(linearises, discards the curvature at projection), $M_B$ projects curvature-faithfully (keeps
the self-similar texture). The load-bearing axis from the previous section --- curvature-fidelity,
not reach --- is thus more precisely \emph{projection-fidelity}: which projection the state
passes through, not what it contains.

\emph{Booking.} The cumulativity of $D(k)$ (the past always present) and the location of the
linearisation in the projection (second level, Doc.~257/301) are derivable and corpus-anchored
respectively; $e$ as a projection error is the more precise form of the earlier formula. The
separation presence / sharpness / linearisation is conceptual bookkeeping.

\section{Inverse inference and the holographic limit}

The model is unrolled from the \emph{measurement side}, for what stands at the centre is the
application: one wants to infer from measurement results back to \emph{causes}. This is the
inverse direction of the reversibility test (Doc.~301) --- not from cause to effect, but from
the measured effect back to the cause. It is exactly this aim that makes the preceding analysis
practically relevant: it asks \emph{how far} a measurement result supports the inference to the
whole.

And here is the limit, which was drawn too weakly before. From a small excerpt one may
\emph{not} infer the whole. The excerpt gives at best a \emph{blurred} image of the whole ---
like a fragment of a holographic recording, which carries the entire image but only unsharply,
at lower resolution. \enquote{The past is present} thus means, more precisely: it is folded in
as a \emph{blurred signature} --- the cumulated value $D(k)$ is a sum from which the summands
cannot be recovered ---, not as a reconstructible recording. \emph{Presence is not
reconstructibility.}

The holographic structure is the self-similarity itself: because every excerpt is
\emph{similar} to the whole, it carries a scaled, blurred image of the total structure. From
this follows something sharper than the earlier talk of \enquote{accuracy}: the sharpness of
the reconstructed whole depends on the \emph{size} of the excerpt (how much of the spiral), not
on the fineness with which a small piece is read. A small window, however losslessly --- down
to the winding quantum --- it is measured, gives a blurred whole; only more excerpt sharpens
it. This is why the absolute reference is missing even at perfect resolution.

The impermissible step --- from part to whole --- is thereby nameable: an \emph{extrapolation
reflex}, taking the fragment for the image. It joins the family of the word, memory, and
linearisation reflexes: in each, something derived, reduced, or partial is taken for the
original, the whole.

For the inverse inference the exact reach follows: inferring from measurement results to causes
is possible, but \emph{only up to the holographic blur}. The rolled-up (compact, 75-exact) side
--- the underlying geometry, the \enquote{cause} --- is \emph{inferred}, not measured; the
unrolled (spiral, measurement) side is where one stands. Rolled and unrolled are here two
readings of the one winding; \enquote{cause} means the direction of the inference, not an
ontological layer beneath the fundamental $D(k)$. Absolute position and ultimate origin
are never sharply recoverable from an excerpt. This connects to the lossiness of accumulation: a
$D(k)$-carrying memory ($M_B$) is a \emph{lossy} accumulator precisely for this reason --- the
holographic blur is the compression.

\emph{Booking.} The measurement-side orientation and the inverse aim (measurement $\to$ cause)
are methodological framing. The holographic limit on inverse inference follows from
self-similarity and reference-lessness (derivable), as does presence $\neq$ reconstructibility.
\enquote{Holographic} is here explicitly an \emph{analogy}: whether every excerpt carries
information about the \emph{whole} (including the distant) --- strong holography --- or merely
resembles its local neighbourhood --- weak self-similarity --- is an open question to be tested
(P35), not an established theorem.

\section{Ordering of the levels}

Three statements must be kept apart, each at its own temperature. \emph{Mathematical}: U3
presupposes U1 and U2, and $D(k)$ is the accumulation structure derivable from the
$K_{\text{frak}}$ recursion (Doc.~295/296); and the linearisation error $e=x-\ln(1+x)$ of a
measurement window is verified against the recursion --- consequences, no speculation marker.
\emph{Philosophical (P35)}: the memory reflex and the linearisation reflex
(measured uniformity as an artefact) as cultural-cognitive inversions, counterparts of the
word reflex (Doc.~301). \emph{Posit}: that time \emph{is} fundamental
(the third level, $D(k)$ as something real that is carried) is the same realist posit as in
Doc.~301, not a new claim. The deflationary reader may always say: all only description. The
argument of this document is independent of that posit --- read even as mere description, U3
remains incoherent without U1 and U2, and the accumulation boundary remains a property of
the time structure, not of memory.

\section{Status}

This document is a statement of position, not a new derivation: it brings together results
from Doc.~285 (time as the woven mass-circle), Doc.~295 (log-spiral, $D(k)$, accumulation
against a bounded kernel), Doc.~296 ($D(k)$ as measurable time, boundarylessness), and
Doc.~301 (three levels, word reflex, self-projection) and Doc.~257 (the winding quantum as
the bit limit) onto the question \enquote{memory or
time first?}. A measurement-theoretic extension places the reference-less excerpt, the congruence/self-similarity
distinction (rolled up: 75 exact parts; unrolled: only asymptotically self-similar) and the
linearisation reflex; the error quantity $e=x-\ln(1+x)$, with the scale-invariant constant
$1-\ln 2$, is verified against the recursion. It is framed as a counter-position to a \enquote{memory $\to$ experienced
time} hierarchy, in a neutral tone and without personal address; the $M_\tau/M_B$ line is
placed, factually, as a cognitive instance of the $D(k)$ structure (booked as a
\emph{bridge candidate}, \emph{restricted}, Category~B), not as its origin. The relationship to the Spiral-Time
Memory Hypothesis is symmetric, the bridge optional and rejectable; provenance and limit are
set out in a dedicated section. The
mathematical statements are corpus-supported; the reflex diagnosis is philosophically
speculative (P35) and booked as such.
